% v8 supplementary section. Existing bibliography: EM.
% Insert after quadric_periods.tex; uses quad:framing and per:chart-definition.
\subsection{二直線と conic の \texorpdfstring{$K_4$}{K4} core}
\label{mixper:section}

五次の分割 $(1,1,1,2)$ に現れる plane core の境界は，四直線ではなく，
二本の直線と一本の conic である．この場合にも，非零 critical periods を選ぶことで
個別の core に作用座標を与えられる．

\begin{proposition}[Mixed plane core の作用座標]\label{mixper:realization}
$L_X,L_Y\subset\PP^2$ を相異なる直線，$C\subset\PP^2$ を smooth conic とし，
$L_X\cup L_Y\cup C$ が三重交点のない SNC divisor であるとする．
正の weights $u_X,u_Y,u_C$ と Kähler form
\[
 [\omega]=2\pi(u_X[L_X]+u_Y[L_Y]+u_C[C])
\]
を固定する．$C$ の crossed quartet と，それぞれの直線の三葉 graph に対応する
marked relative isotopy classes を指定する．
その中で代表を選べば，定理~\ref{snc:main} による core に，
定義~\ref{per:chart-definition} の作用座標に適合する標準 $K_4$ cone chart を構成できる．
六つの face maps は unimodular で，branch denominators は全て一である．
\end{proposition}

\begin{proof}
link の頂点を，二直線に対応する $X,Y$ と，conic の二 saddles に対応する $U,V$ とする．
五つの crossing edges は $XY,XU,XV,YU,YV$，内部 annulus edge は $UV$ である．
crossed quartet なので，$U,V$ はそれぞれ $L_X\cap C,L_Y\cap C$ の一点ずつに接する．
標準格子は $N_4=\ZZ^4/\ZZ(1,1,1,1)$，rays は
$r_X,r_Y,r_U,r_V$ とする．

$W=u_X+u_Y+2u_C$ と置くと，二直線の面積は $2\pi W$，
conic の面積は $4\pi W$ である．crossing germs の準備後に foliation の代表を選ぶ．
各 puncture basin の面積を $2\pi$ で割った値を次のように指定する．
\[
\begin{array}{c|ccc}
 \text{component}&\multicolumn{3}{c}{\text{basin targets}}\\\hline
 L_X&Y:u_Y+u_X&U:u_C&V:u_C\\
 L_Y&X:u_X&U:u_C+u_Y&V:u_C
\end{array}
\]
conic では
\[
 U\to X:u_X+u_C,\quad U\to Y:u_Y,\quad
 V\to X:u_X+u_C,\quad V\to Y:u_Y+u_C
\]
とする．二直線では三つの target の和が $W$ であり，conic では四つの和が
$2u_X+2u_Y+3u_C=2W-u_C$ となる．従って conic の内部 annulus の面積は $2\pi u_C$ である．
全ての領域の面積は正なので，三 basin および四 basin の relative Moser の議論により，
十分小さい crossing disks を固定して実現できる．

これを有向 crossing graph で表すと
\begin{equation}\label{mixper:orientation}
 X\longrightarrow Y,\quad Y\longrightarrow U,\quad
 U\longrightarrow X,\quad V\longrightarrow X,\quad V\longrightarrow Y
\end{equation}
となる．矢印の始点側の critical slope が $-1$，終点側が零である．
これは一般式 $h=(u_{\mathrm{other}}-\operatorname{area}(D)/(2\pi))/u_{\mathrm{own}}$ による．
graph-period correction は fiber tangent を消す basic form によるため，この作用計算を変えない．
各 vertex $v$ の common normal meridian を $\alpha_v$ とすれば，
branch ray parameter を $s=-I_{\alpha_v}>0$ と取れる．

crossing face $vw$ の基底を，対応する二成分の正の meridians $(\mu_v,\mu_w)$ とする．
$\epsilon_{vw}$ を $v\to w$ が選ばれた場合に一，それ以外に零と置いて
\[
 \begin{pmatrix}I_{\mu_v}\\I_{\mu_w}\end{pmatrix}
 =\begin{pmatrix}-1&\epsilon_{wv}\\\epsilon_{vw}&-1\end{pmatrix}
   \begin{pmatrix}x\\y\end{pmatrix},\qquad
 F=xr_v+yr_w
\]
とする．\eqref{mixper:orientation} では各 crossing にちょうど一つの向きがあるため，
五つの determinant は全て一である．

conic の normal meridian を $\alpha_C$，$U$ 側の二つの puncture cycles の和を
内部 torus に運んだものを $\beta_U$，$V$ 側のものを $\beta_V$ とする．
normal Euler 数は $C^2=4$ なので，補題~\ref{quad:framing} の clutching 計算で二を四に置き換えると，
実際の annular torus 上で
\[
 \beta_U=-4\alpha_C-\beta_V
\]
となる．$(\alpha_C,\beta_U)$ は整数基底である．
$U$ 側の二つの critical slopes の和は $-1$，$V$ 側は $-2$ なので，
この基底における両端の action directions は $(-1,1)$ と $(-1,2)$ である．
従って内部 face の map は
\[
 \begin{pmatrix}I_{\alpha_C}\\I_{\beta_U}\end{pmatrix}
 =\begin{pmatrix}-1&-1\\1&2\end{pmatrix}
   \begin{pmatrix}x\\y\end{pmatrix},\qquad
 F=xr_U+yr_V
\]
で与えられ，determinant は $-1$ である．

正値性は全 face で確認できる．crossing の選択された側の arm では partial slope が
$h>-1$，反対側では $h>0$ なので，inverse matrix の二係数は正である．
transition caps と product crossing では二つの meridian actions を負に保つように normal radius を縮める．
conic の内部 annulus では $I_{\beta_U}/I_{\alpha_C}$ が $-1$ から $-2$ へ厳密単調に変わり，
$x=-2I_{\alpha_C}-I_{\beta_U}>0$，$y=I_{\alpha_C}+I_{\beta_U}>0$ となる．

三つの incident faces は各 vertex $v$ に沿って $F=-I_{\alpha_v}r_v$ に一致する．
全 endpoint directions の normal component が $-1$ なので，branch lattice は飽和している．
標準 $K_4$ fan の全 cones はそれぞれの rank-two span で unimodular であり，
上の六つの action matrices も unimodular である．
Liouville homogeneity と contact-link 上の projective action ratio の厳密単調性により，
各 face の map は所期の cone sector を一度だけ被覆する\cite[Remarks 2.11--2.12]{EM}．
compact link を用いて core を含む近傍へ貼り合わせれば，主張した chart を得る．
\end{proof}

\begin{remark}[四つの最小整数 recipes]\label{mixper:four}
各直線で一つの basin をその直線の weight 分だけ拡大し，conic では四 basins 中三つを
それぞれ $u_C$ 分だけ拡大する方法には $3^2\cdot4=36$ 個の marked choices がある．
そのうち全 face を unimodular にするものは正確に四つである．
実際，選んだ half-edges の総数が crossing edges の数五に等しいので，
重複を避けるには全 crossing をちょうど一回選ぶ必要がある．
$XY$ を $X\to Y$ と向けた場合，$X$ の唯一の選択はそれで尽くされ，
$U\to X,V\to X$ が強制される．$Y$ は $Y\to U$ または $Y\to V$ の二択となる．
$Y\to X$ の場合も二択で，合計四つである．
内部 action matrix は $\left(\begin{smallmatrix}-1&-1\\m&m+1\end{smallmatrix}\right)$，
$m\in\{1,2\}$ となるので，全て同じ証明で chart を与える．
この有限個数は限定した面積 recipe を数えている．
\end{remark}

\begin{corollary}[Mixed quintic の各 core の個別実現]\label{mixper:individual}
五次の分割 $(1,1,1,2)$ で全 core links を matroidal にする四つの標識付き quartet systems の各々について，
一つの core ずつ考えれば，全ての core が作用座標に適合する標準 fan chart を持つように，
marked foliations の代表を選べる．
三つの plane cores は $K_4$，quadric core は $K_{3,3}$ である．
\end{corollary}

\begin{proof}
三つの plane cores に命題~\ref{mixper:realization}，quadric core に定理~\ref{quad:realization} を適用する．
\end{proof}

この系は，異なる core のために選んだ代表や面積指定を共有 actual curves 上で同時に一致させるとは述べていない．
既存の branch graph と標準 coarse cones を保った大域 atlas への接続は，
後の定理~\ref{hd9:mixed-no-go} により任意の denominators に対して不可能である．
branch graph や core，cone の格子を変更する場合は別の問題として残る．
この大域障害にも，本節の個別構成にも，全 branch blocks の周期を
ambient linear forms に延長する条件は課していない．
